\section{Random dense recursive inner (log-memory baseline)}
\label{sec:inner-dense}

\paragraph{Role in the modular serial concatenation framework.}
In Section~\ref{sec:framework} (the modular serial concatenation framework),
we reduce minimum-distance analysis of the concatenated code
\[
\Csc \ :=\ \Ctwo \circ \Pi \circ \Cone \ \subseteq\ \{0,1\}^N
\]
to bounding, for each outer weight $w$, the inner tail probability
\[
p_w(\delta)\ :=
\PP\!\left[\ \wt\!\big(\Enc_{\mathrm{in}}(U)\big)\le \delta N \ \middle|\ \wt(U)=w\ \right],
\]
where $U$ is uniform over $\{0,1\}^N$ conditioned on $\wt(U)=w$.
(Equivalently: $\Pi$ is uniform, so $\Pi(\text{any fixed weight-}w\text{ word})$ is uniform on the weight-$w$ slice.)

This section instantiates the inner code $\Ctwo$ as a \emph{time-varying fixed-tap dense recursive
encoder} with memory $m=\Theta(\log N)$, and proves two modular bounds on $p_w(\delta)$:
\begin{enumerate}
\item a \textbf{general three-term tail bound} (useful for large $w$ / large deviations), and
\item a \textbf{forced-termination contraction bound} that converts the ``OFF mechanism'' into a \emph{geometric-in-$w$}
term for small/intermediate $w$.
\end{enumerate}
These are the only inner facts later sections need.

% ============================================================
\subsection{Inner ensemble definition (fixed-tap dense recursive, rate 1)}
\label{subsec:inner-dense-def}

Fix blocklength $N$ and a memory parameter
\[
m \ :=\ \gamma \log_2 N
\qquad\text{for a constant }\gamma>0.
\]
The inner encoder is a streaming linear map $\Enc_{\mathrm{in}}:\F_2^N\to\F_2^N$ defined via a shift-register state.
The construction used in the dense+dense numerics is the fixed-tap variant: the oldest state tap is fixed to one,
while the remaining taps are fresh dense random bits.

\paragraph{Streaming definition.}
Initialize $s_1:=0^m\in\F_2^m$.
For each time $t\in[N]$:
\begin{enumerate}
\item Sample a fresh random vector \(a_t\gets\Unif(\F_2^{m-1})\), independently over \(t\), and set
\[
k_t := (1,a_t)\in \F_2^m.
\]
\item Output
\begin{equation}
\label{eq:dense-inner-def}
y_t \ :=\ u_t \oplus \langle k_t, s_t\rangle,
\qquad
s_{t+1}\ :=\ \Shift(s_t,y_t),
\end{equation}
where $\Shift$ drops the oldest state bit and appends $y_t$.
\end{enumerate}
Equivalently, for $t\ge 2$, the state stores the last $m$ outputs
\[
s_t \ =\ (y_{t-m},y_{t-m+1},\dots,y_{t-1})
\quad\text{with zero padding for indices }\le 0.
\]

\paragraph{Linearity and generator matrix.}
For a fixed draw of $(k_t)_{t=1}^N$, the map $u\mapsto y$ is $\F_2$-linear and systematic ($y_t$ contains $u_t$),
hence it corresponds to multiplication by an invertible $N\times N$ binary matrix $\Gtwo$:
\[
y \ =\ u \cdot \Gtwo.
\]
We view $\Ctwo$ as the (rate-1) code induced by this random $\Gtwo$ ensemble.

\paragraph{ON/OFF notation.}
Define events
\[
\ON_t \ :=\ \{s_t\neq 0\},
\qquad
\OFF_t\ :=\ \{s_t=0\}.
\]
Also define the critical termination-ready state
\[
\CRIT_t \ :=\ \{s_t=(1,0,\dots,0)\},
\qquad
\FREE_t \ :=\ \ON_t\cap \CRIT_t^c.
\]
Intuitively: in \OFF, the encoder outputs $y_t=u_t$ deterministically. In \(\FREE\), at least one random tap is exposed,
so the output is a fresh unbiased bit. In \(\CRIT\), the fixed leading tap makes the output deterministic:
\[
y_t = u_t\oplus 1.
\]
Thus a fixed-tap ON episode can terminate only when the input has a \(1\) exactly at a critical time.

% ============================================================
\subsection{Core ON/OFF lemmas}
\label{subsec:inner-dense-onoff}

Let $\mathcal{F}_{t-1}:=\sigma(u_1^t,k_1^{t-1},y_1^{t-1})$ be the natural past up to time $t-1$.

\begin{lemma}[Fresh fair coin while free]
\label{lem:dense-fair-while-on}
For every $t$, conditioned on $\mathcal{F}_{t-1}$ and on $\FREE_t$, the output bit $y_t$ is uniform:
\[
\PP[y_t=0\mid \mathcal{F}_{t-1},\FREE_t]=\PP[y_t=1\mid \mathcal{F}_{t-1},\FREE_t]=\tfrac12.
\]
\end{lemma}

\begin{proof}
Fix $t$. On $\FREE_t$, the state \(s_t\) is nonzero and is not \((1,0,\dots,0)\). Hence at least one of the last
\(m-1\) state coordinates is nonzero. Since the corresponding coordinates of \(a_t\) are uniform and independent of
\(\mathcal{F}_{t-1}\), the random part of \(\langle k_t,s_t\rangle\) is a nontrivial linear functional of a uniform
vector, hence uniform in \(\F_2\). Therefore \(y_t=u_t\oplus\langle k_t,s_t\rangle\) is uniform.
\end{proof}

\begin{lemma}[Free-time fair domination]
\label{lem:dense-on-time-domination}
Fix a history up to time \(a-1\). Let
\(\tau_1<\tau_2<\cdots\) be any sequence of predictable stopping times, adapted to the encoder history before the
current output is revealed, such that on the event \(\{\tau_j<\infty\}\) the encoder is in \(\FREE\) at time
\(\tau_j\). Then
for every \(M\ge 1\) and every \(B\subseteq\{0,1\}^M\),
\[
\PP\!\left[
\tau_M<\infty\ \wedge\ (Y_{\tau_1},\dots,Y_{\tau_M})\in B
\mid \mathcal{F}_{a-1}
\right]
\le
|B|\,2^{-M}.
\]
In particular,
\[
\PP\!\left[
\tau_M<\infty\ \wedge\ \sum_{j=1}^M Y_{\tau_j}\le d
\mid \mathcal{F}_{a-1}
\right]
\le
\PP[\Binomial(M,\tfrac12)\le d].
\]
\end{lemma}

\begin{proof}
It is enough to prove the first display for a fixed string \(b\in\{0,1\}^M\), then sum over \(b\in B\).
Reveal the process up to the successive stopping times. On the event that \(\tau_j\) is finite, \(\FREE_{\tau_j}\)
holds by assumption, and Lemma~\ref{lem:dense-fair-while-on} gives
\[
\PP[Y_{\tau_j}=b_j\mid \mathcal{F}_{\tau_j-1},\FREE_{\tau_j}]=\tfrac12.
\]
The fresh tap vector at time \(\tau_j\) is independent of the past, so multiplying these conditional bounds over
\(j=1,\dots,M\) gives \(2^{-M}\). The binomial-tail consequence follows by taking
\(B=\{b:\sum_j b_j\le d\}\).
\end{proof}

\begin{lemma}[Surviving fixed-tap episodes dominate fair coins]
\label{lem:fixedtap-surviving-episode-tail}
Fix an input word and suppose an ON episode has started. Let \(E_L\) be the event that this episode survives for at
least its next \(L\) output times. Then for every \(d\),
\[
\PP\!\left[
E_L\ \wedge\ \sum_{j=1}^L Y_{\tau+j}\le d
\right]
\le
\PP[\Binomial(L,\tfrac12)\le d],
\]
where \(\tau+1,\dots,\tau+L\) are the first \(L\) output times after the episode starts.
\end{lemma}

\begin{proof}
Expose the \(L\) output times sequentially and condition on the past together with survival up to the current time.
If the state is in \(\FREE\), Lemma~\ref{lem:dense-fair-while-on} gives a fresh fair bit. If the state is in
\(\CRIT\), then survival through the current time rules out the aligned terminating case \(u_t=1,y_t=0\), and the
fixed tap forces \(y_t=1\). Thus, at every step on the surviving path, the next output has conditional probability of
being \(1\) at least \(1/2\). The usual induction for Bernoulli stochastic domination gives the displayed lower-tail
bound.
\end{proof}

\begin{lemma}[Exact final-survivor suffix average]
\label{lem:fixedtap-final-survivor-suffix}
Let \(U\) be uniform on the weight-\(h\) slice of \(\{0,1\}^N\), and let
\[
\tau(U):=\min\supp(U),\qquad S(U):=N-\tau(U)+1.
\]
Let \(\mathcal{S}_{\mathrm{fin}}\) be the event that the ON episode started by the first input one survives through the
end of the block. Then, for every \(d\ge 0\),
\[
\PP\!\left[
\mathcal{S}_{\mathrm{fin}}\wedge \wt(Y)\le d
\middle| \wt(U)=h
\right]
\le
\sum_{S=h}^{N}
\frac{\binom{S-1}{h-1}}{\binom{N}{h}}\,
\PP[\Binomial(S,\tfrac12)\le d].
\]
\end{lemma}

\begin{proof}
Condition on the first input-one position \(\tau=t\), and put \(S=N-t+1\). The bit \(y_t\) is deterministically \(1\),
because the encoder is still \(\OFF\) before the first input one. On \(\mathcal{S}_{\mathrm{fin}}\), the remaining
outputs \(y_{t+1},\dots,y_N\) stay on a surviving ON episode. By the same stochastic-domination argument as
Lemma~\ref{lem:fixedtap-surviving-episode-tail}, the suffix output weight
\(\sum_{i=t}^{N}Y_i\) dominates \(\Binomial(S,\tfrac12)\) in the lower-tail direction. Hence, for this fixed value of
\(t\),
\[
\PP[\mathcal{S}_{\mathrm{fin}}\wedge \wt(Y)\le d\mid U,\tau=t]
\le
\PP[\Binomial(S,\tfrac12)\le d].
\]
Finally, under the uniform weight-\(h\) law,
\[
\PP[S(U)=S\mid \wt(U)=h]
=
\frac{\binom{S-1}{h-1}}{\binom{N}{h}},
\]
since the first one is at position \(N-S+1\) and the remaining \(h-1\) ones are chosen from the later \(S-1\) positions.
Averaging over \(S\) proves the claim.
\end{proof}

\begin{lemma}[ON$\to$OFF requires an $m$-run of zeros in the output]
\label{lem:dense-off-requires-mzeros}
If $\OFF_q$ holds at some time $q\ge 2$, then the preceding $m$ outputs are all zero:
\[
y_{q-m}=y_{q-m+1}=\cdots=y_{q-1}=0.
\]
Equivalently, every ON$\to$OFF transition is witnessed by a contiguous length-$m$ all-zero window in the output stream.
\end{lemma}

\begin{proof}
By construction, $s_q$ is the vector of the last $m$ outputs before time $q$. Thus $s_q=0$ iff those $m$ outputs are all zero.
\end{proof}

\begin{lemma}[Fixed-tap termination requires an aligned input one]
\label{lem:fixedtap-termination-needs-input-one}
Suppose the fixed-tap dense inner makes an ON\(\to\)OFF transition at time \(q\), meaning \(\ON_{q-1}\) and
\(\OFF_q\). Then
\[
s_{q-1}=(1,0,\dots,0),
\qquad
u_{q-1}=1,
\]
and the output contains the terminating pattern
\[
y_{q-m-1}=1,
\qquad
y_{q-m}=y_{q-m+1}=\cdots=y_{q-1}=0.
\]
Equivalently, every genuine fixed-tap termination is a \(10^m\) output pattern whose last zero is aligned with an input
one.
\end{lemma}

\begin{proof}
Since \(\OFF_q\), Lemma~\ref{lem:dense-off-requires-mzeros} gives
\(y_{q-m}=\cdots=y_{q-1}=0\). Since \(\ON_{q-1}\), the previous state
\[
s_{q-1}=(y_{q-m-1},y_{q-m},\dots,y_{q-2})
\]
is nonzero. The last \(m-1\) displayed coordinates are zero, so the first coordinate must be one; hence
\(s_{q-1}=(1,0,\dots,0)\). In the fixed-tap encoder the leading tap is fixed to one and all other state coordinates are
zero at this time, so
\[
y_{q-1}=u_{q-1}\oplus 1.
\]
But \(y_{q-1}=0\), hence \(u_{q-1}=1\).
\end{proof}

\begin{lemma}[Fixed-tap termination candidates]
\label{lem:fixedtap-termination-candidates}
Fix an input word \(u\). For any interval \(I\subseteq[N]\), the probability that an ON\(\to\)OFF transition has its
last output time in \(I\) is at most
\[
|\{t\in I:u_t=1\}|\,2^{-(m-1)}.
\]
\end{lemma}

\begin{proof}
By Lemma~\ref{lem:fixedtap-termination-needs-input-one}, a transition with last output time \(t\) requires \(u_t=1\)
and \(s_t=(1,0,\dots,0)\). For a fixed such \(t\), the state condition requires the preceding \(m-1\) outputs after
the leading \(1\) in the state to be zero. None of these zero outputs can be produced from the critical state without
having already terminated earlier: if the input is zero, the critical state forces output \(1\); if the input is one,
the critical state would terminate. Hence those \(m-1\) zero outputs occur at free times. Applying
Lemma~\ref{lem:dense-on-time-domination} to those predictable free times gives probability at most \(2^{-(m-1)}\).
Union bound over the input-one positions in \(I\).
\end{proof}

\begin{lemma}[Fixed-tap episode budget]
\label{lem:fixedtap-episode-budget}
For a fixed input-output trajectory of the fixed-tap dense recursive inner, let
\[
S:=|\{t:\OFF_t,\ \ON_{t+1}\}|,
\qquad
D:=|\{t:\ON_t,\ \OFF_{t+1}\}|
\]
be the number of genuine starts and genuine terminations inside the block. Then
\[
S+D \le \wt(U).
\]
Moreover, each termination time is an input-one position that is not a start time.
\end{lemma}

\begin{proof}
If \(\OFF_t\) and \(\ON_{t+1}\), then \(s_t=0\), so \(y_t=u_t\). Therefore a genuine start requires \(u_t=1\).
If \(\ON_t\) and \(\OFF_{t+1}\), Lemma~\ref{lem:fixedtap-termination-needs-input-one} gives \(u_t=1\) as well.
The two cases are disjoint at a fixed time \(t\), since \(s_t\) cannot be both zero and nonzero. Thus starts and
terminations inject into disjoint subsets of the support of \(U\), proving \(S+D\le \wt(U)\).
\end{proof}

\begin{lemma}[One fixed-tap episode can terminate only at later input ones]
\label{lem:fixedtap-one-episode-termination}
Fix an input word \(u\), and suppose an ON episode starts immediately after an input-one position \(i\). For
\(L\ge 1\), let
\[
B_i(u;L):=
|\{j\in \supp(u): i<j\le i+L\}|.
\]
Then the probability that this episode terminates within its first \(L\) output times is at most
\[
B_i(u;L)\,2^{-(m-1)}.
\]
\end{lemma}

\begin{proof}
By Lemma~\ref{lem:fixedtap-termination-needs-input-one}, a termination whose last output time is in
\(\{i+1,\dots,i+L\}\) must be aligned with an input-one position in that interval. There are \(B_i(u;L)\) such
positions. For each fixed candidate, Lemma~\ref{lem:fixedtap-termination-candidates} gives probability at most
\(2^{-(m-1)}\), and the claim follows by a union bound.
\end{proof}

\begin{lemma}[Fixed-tap multi-candidate termination bound]
\label{lem:fixedtap-multicandidate-termination}
Fix an input word \(u\), and fix \(r\) distinct input-one positions
\[
t_1<\cdots<t_r.
\]
The probability that all \(r\) of these positions are genuine fixed-tap termination times is at most
\[
2^{-r(m-1)}.
\]
\end{lemma}

\begin{proof}
Expose the trajectory from left to right. At a genuine fixed-tap termination time \(t_a\),
Lemma~\ref{lem:fixedtap-termination-needs-input-one} forces the state just before output \(t_a\) to be
\((1,0,\dots,0)\). Equivalently, after the leading \(1\) in the state, the next \(m-1\) relevant outputs must be zero
while the encoder is not already terminated. These are free-time zero constraints in the sense of
Lemma~\ref{lem:dense-on-time-domination}, except on histories where the requested termination pattern is already
impossible. Thus, conditional on any past compatible with the earlier requested terminations, the next requested
termination costs at most \(2^{-(m-1)}\). Multiplying over \(a=1,\dots,r\) gives the claim.
\end{proof}

\begin{corollary}[Isolated fixed-tap termination candidates are pair-controlled]
\label{cor:fixedtap-isolated-pair-candidates}
Let \(U\) be uniform on the isolated slice \(\{\wt(U)=h,\ R(U)=h\}\). Let
\[
\mathcal{P}_L(U)
:=
|\{(i,j): i<j,\ U_i=U_j=1,\ j-i\le L\}|
\]
be the number of isolated input-one pairs with separation at most \(L\). Then
\[
\EE[\mathcal{P}_L(U)\mid \wt(U)=h,\ R(U)=h]
\le
\binom{h}{2}\frac{2L}{N-h}.
\]
Consequently, the expected number of possible within-\(L\) fixed-tap termination candidates over all isolated starts is
bounded by the same quantity.
\end{corollary}

\begin{proof}
On the isolated slice, write the support positions as
\[
1\le x_1<\cdots<x_h\le N,
\qquad x_{a+1}-x_a\ge 2.
\]
The shifted coordinates \(z_a:=x_a-(a-1)\) form a uniformly random \(h\)-subset of \([M]\), where \(M=N-h+1\). For any
unordered pair of selected shifted coordinates, the probability that their distance is at most \(L\) is at most
\[
\frac{ML}{\binom{M}{2}}\le \frac{2L}{M-1}=\frac{2L}{N-h},
\]
because there are at most \(ML\) unordered pairs in \([M]\) with separation at most \(L\). Since
\(x_b-x_a=(z_b-z_a)+(b-a)\), the event \(x_b-x_a\le L\) is contained in \(z_b-z_a\le L\). Linearity of expectation
then proves the displayed bound.
\end{proof}

% ============================================================
\subsection{A general three-term tail bound for $p_w(\delta)$}
\label{subsec:inner-dense-three-term}

We first prove a bound that is convenient for large $w$ and for the linear-weight regime:
it isolates three distinct mechanisms by which $\wt(Y)$ can be small.

\begin{lemma}[First 1 occurs early for uniform weight-$w$ inputs]
\label{lem:first-one-early}
Let $U$ be uniform over $\{0,1\}^N$ conditioned on $\wt(U)=w\ge 1$.
Let $T:=\min\{t:U_t=1\}$ be the first index containing a $1$.
Then for any $\varepsilon\in(0,1)$,
\[
\PP[T>\varepsilon N]\ \le\ (1-\varepsilon)^w.
\]
\end{lemma}

\begin{proof}
$T>\varepsilon N$ means the first $\lfloor \varepsilon N\rfloor$ coordinates contain no $1$.
Under the uniform weight-$w$ model,
\[
\PP[T>\varepsilon N]
=
\frac{\binom{N-\lfloor\varepsilon N\rfloor}{w}}{\binom{N}{w}}
=
\prod_{i=0}^{w-1}\frac{N-\lfloor\varepsilon N\rfloor-i}{N-i}
\le
\left(\frac{N-\lfloor\varepsilon N\rfloor}{N}\right)^w
\le
(1-\varepsilon)^w.
\]
\end{proof}

\begin{corollary}[Zero-input gaps survive in the fixed-tap inner]
\label{cor:dense-gap-survival}
Fix an interval of $g\ge 0$ consecutive zero inputs, and condition on the event that the encoder is ON at the start of
that interval. Then the encoder remains ON throughout the whole gap:
\begin{equation}
\label{eq:dense-gap-survival}
\PP[\text{survive a zero gap of length }g \mid \ON]
\;=\;
1,
\end{equation}
apart from the harmless convention at the block boundary after time \(N\).
\end{corollary}

\begin{proof}
By Lemma~\ref{lem:fixedtap-termination-needs-input-one}, a genuine ON\(\to\)OFF transition requires the final zero of
the terminating \(10^m\) pattern to be aligned with an input one. Inside a zero-input gap no such aligned one exists,
so no ON\(\to\)OFF transition can occur there.
\end{proof}

\begin{lemma}[Dense inner: three-term low-weight tail]
\label{lem:dense-three-term-tail}
Let $U$ be uniform over $\{0,1\}^N$ conditioned on $\wt(U)=w\ge 1$, and let $Y=\Enc_{\mathrm{in}}(U)$ be produced by
\eqref{eq:dense-inner-def}. Fix $\delta\in(0,1/2)$ and $\varepsilon\in(0,1)$. Then
\begin{equation}
\label{eq:dense-three-term}
\PP[\wt(Y)\le \delta N \mid \wt(U)=w]
\ \le\
(1-\varepsilon)^w\ +\ N\cdot 2^{-m}\ +\ q_{\delta,\varepsilon}(N),
\end{equation}
where
\[
q_{\delta,\varepsilon}(N)
:=
\PP\!\left[\ \mathrm{Binomial}(\lfloor(1-\varepsilon)N\rfloor,\tfrac12)\ \le\ \delta N\ \right]
\ \le\
2^{-(1-\varepsilon)N\cdot \big(1-H_2(\delta/(1-\varepsilon))\big)}
\]
whenever $\delta\le (1-\varepsilon)/2$.
\end{lemma}

\begin{proof}
Let $T$ be the first index with $U_T=1$. By Lemma~\ref{lem:first-one-early},
\[
\PP[T>\varepsilon N]\le (1-\varepsilon)^w.
\]
On the event $\{T\le \varepsilon N\}$ we have $s_T=0$ (since $s_1=0$ and no earlier $1$ has occurred),
hence $Y_T=U_T=1$, which implies $\ON_{T+1}$.

If the encoder ever transitions from \ON to \OFF after time $T+1$, then by Lemma~\ref{lem:dense-off-requires-mzeros}
there exists a length-$m$ all-zero window in $Y$ starting at some $r\in\{T+1,\dots,N-m+1\}$.
For a fixed $r$, Lemma~\ref{lem:dense-on-time-domination} gives
\[
\PP[Y_r=\cdots=Y_{r+m-1}=0 \ \wedge\ \ON_r]\le 2^{-m}.
\]
A union bound over all $r$ yields
\[
\PP[\exists\ \text{such window}] \ \le\ N\cdot 2^{-m}.
\]

On the complementary event that $T\le \varepsilon N$ and no ON$\to$OFF transition occurs after $T+1$,
we have $\ON_t$ for all $t\in\{T+1,\dots,N\}$, so Lemma~\ref{lem:dense-on-time-domination} bounds the joint
low-weight probability by the corresponding fair-coin lower tail.
Since $T\le \varepsilon N$ implies $N-T\ge (1-\varepsilon)N$, monotonicity of binomial tails yields the term
$q_{\delta,\varepsilon}(N)$.

Union bounding these three mechanisms proves \eqref{eq:dense-three-term}.
\end{proof}

\paragraph{Interpretation.}
The three terms in \eqref{eq:dense-three-term} correspond to:
(i) the first $1$ appears too late (probability $(1-\varepsilon)^w$),
(ii) an ON$\to$OFF transition occurs (probability $\le N2^{-m}$),
(iii) otherwise, the encoder stays ON for at least $(1-\varepsilon)N$ steps and outputs behave like fair coins, so
$\wt(Y)$ has a large-deviation tail $q_{\delta,\varepsilon}(N)$.

\begin{lemma}[First-start-aware no-OFF tail]
\label{lem:dense-first-start-nooff-tail}
Let \(U\) be uniform over \(\{0,1\}^N\) conditioned on \(\wt(U)=w\ge 1\), and let
\[
T:=\min\{t:U_t=1\}.
\]
Let \(Y=\Enc_{\mathrm{in}}(U)\), and let \(\mathcal{E}_{\mathrm{off}}\) be the event that after the first activation
at time \(T\), the encoder later returns to \(\OFF\). Then for every integer \(d\ge 0\),
\begin{equation}
\label{eq:dense-first-start-nooff-tail}
\PP[\wt(Y)\le d\ \wedge\ \mathcal{E}_{\mathrm{off}}^c\mid \wt(U)=w]
\;\le\;
\sum_{t=1}^{N-w+1}
\frac{\binom{N-t}{w-1}}{\binom{N}{w}}\,
\PP[\Binomial(N-t,\tfrac12)\le d].
\end{equation}
Consequently,
\begin{equation}
\label{eq:dense-first-start-with-off-remainder}
p_w(\delta)
\;\le\;
\sum_{t=1}^{N-w+1}
\frac{\binom{N-t}{w-1}}{\binom{N}{w}}\,
\PP[\Binomial(N-t,\tfrac12)\le \lfloor\delta N\rfloor]
\;+\;
\PP[\mathcal{E}_{\mathrm{off}}\mid \wt(U)=w].
\end{equation}
\end{lemma}

\begin{proof}
The exact first-start law under the uniform weight-\(w\) model is
\[
\PP[T=t\mid \wt(U)=w]
=
\frac{\binom{N-t}{w-1}}{\binom{N}{w}},
\qquad 1\le t\le N-w+1.
\]
Condition on \(T=t\). At time \(t\), the encoder is still \(\OFF\), so \(Y_t=U_t=1\) and the next state is nonzero.
On \(\mathcal{E}_{\mathrm{off}}^c\), the encoder remains \(\ON\) throughout the suffix \(t+1,\dots,N\). Hence
Lemma~\ref{lem:dense-on-time-domination} bounds the joint probability that the suffix has at most \(d\) ones by
\(\PP[\Binomial(N-t,\tfrac12)\le d]\), which gives \eqref{eq:dense-first-start-nooff-tail} after averaging over \(T\).
The final displayed bound follows by adding the OFF event.
\end{proof}

\paragraph{Position-aware first-start tail.}
For \(w=\eta N\) and \(t=\alpha N\), the summand in
\eqref{eq:dense-first-start-nooff-tail} has exponent
\[
H_2(\eta)
-
(1-\alpha)H_2\!\left(\frac{\eta}{1-\alpha}\right)
\;+\;
(1-\alpha)
\left[
1-H_2\!\left(\frac{\delta}{1-\alpha}\right)
\right],
\]
whenever \(\alpha\le 1-\eta\) and \(\delta/(1-\alpha)<1/2\). Optimizing over \(\alpha\) gives the correct
first-start-aware replacement for the earlier coarse \((1-\varepsilon)^w\) splice. The remaining linear-window task is
to replace the raw OFF remainder in \eqref{eq:dense-first-start-with-off-remainder} by an exponentially small
restart/episode bound.

\begin{corollary}[Single-run base factor]
\label{cor:dense-inner-single-run}
Fix $\delta\in(0,1/2)$ and let $U$ be uniform over $\{0,1\}^N$ conditioned on $\wt(U)=1$.
Let $Y=\Enc_{\mathrm{in}}(U)$ be produced by \eqref{eq:dense-inner-def}. Then
\begin{equation}
\label{eq:dense-inner-single-run}
\PP[\wt(Y)\le \delta N \mid \wt(U)=1]
\;\le\;
\frac{\lceil 2\delta N\rceil}{N}
\;+\;
N2^{-m}
\;+\;
\PP[\Binomial(\lceil 2\delta N\rceil,\tfrac12)\le \delta N].
\end{equation}
In particular,
\begin{equation}
\label{eq:dense-inner-single-run-asym}
\PP[\wt(Y)\le \delta N \mid \wt(U)=1]
\;\le\;
2\delta
\;+\;
O\!\left(\frac{1}{N}\right)
\;+\;
N2^{-m}
\;+\;
2^{-\Omega_\delta(N)}.
\end{equation}
\end{corollary}

\begin{proof}
Let $T$ be the unique index with $U_T=1$. Under the uniform weight-$1$ model, $T$ is uniform on $\{1,\dots,N\}$.

On the event that the encoder never transitions from \ON to \OFF after time $T+1$,
Lemma~\ref{lem:dense-on-time-domination} bounds every low-weight suffix event by the corresponding fair-coin tail.
If $N-T\ge \lceil 2\delta N\rceil$, then
\[
\PP[\wt(Y)\le \delta N \mid T,\ \text{no OFF after }T+1]
\;\le\;
\PP[\Binomial(\lceil 2\delta N\rceil,\tfrac12)\le \delta N]
\]
by monotonicity in the number of fair-coin steps. Otherwise $N-T<\lceil 2\delta N\rceil$, which is equivalent to
\[
T > N-\lceil 2\delta N\rceil.
\]
Since $T$ is uniform, this exceptional late-start event has probability at most $\lceil 2\delta N\rceil/N$.

Finally, by Lemma~\ref{lem:dense-off-requires-mzeros} and the same union bound as in Lemma~\ref{lem:dense-three-term-tail},
the probability that an ON$\to$OFF transition occurs after time $T+1$ is at most $N2^{-m}$.
Union bounding these three mechanisms proves \eqref{eq:dense-inner-single-run}.

For \eqref{eq:dense-inner-single-run-asym}, note that $\lceil 2\delta N\rceil/N = 2\delta + O(1/N)$, while the
binomial term is exponentially small because $\delta < \lceil 2\delta N\rceil/(2\lceil 2\delta N\rceil)$ for all
sufficiently large $N$.
\end{proof}

% ============================================================
\subsection{Forced-termination contraction (tightening the small-$w$ regime)}
\label{subsec:inner-dense-forced-termination}

For small/intermediate weights, the pointwise polynomial term $N2^{-m}$ in Lemma~\ref{lem:dense-three-term-tail}
is undesirable. We now exploit recursion plus the run structure of a sparse interleaved input to turn the OFF mechanism
into a genuine geometric-in-$w$ contraction.

\begin{definition}[Runs of ones]
\label{def:runs}
For $u\in\{0,1\}^N$, let $R(u)$ denote the number of runs of ones in $u$,
i.e., the number of maximal contiguous substrings of $1$'s.
\end{definition}

\begin{definition}[Post-start OFF time]
\label{def:post-start-off-time}
For an input-output trajectory, let \(T=\min\{t:U_t=1\}\), and define
\[
\mathrm{OffAfter}(T)
\;:=\;
\bigl|\{\,t\in\{T+1,\dots,N\}: s_t=0\,\}\bigr|.
\]
\end{definition}

\begin{lemma}[Exact run-count distribution under the uniform weight-$w$ model]
\label{lem:run-count-distribution}
Let $U$ be uniform over $\{0,1\}^N$ conditioned on $\wt(U)=w\ge 1$.
Then for every $s\in\{1,\dots,w\}$,
\[
\PP[R(U)=s]
\;=\;
\frac{\binom{w-1}{s-1}\binom{N-w+1}{s}}{\binom{N}{w}}.
\]
Consequently, for every $r\in\{1,\dots,w\}$,
\begin{equation}
\label{eq:run-lower-tail-exact}
\PP[R(U)<r]
\;=\;
\frac{1}{\binom{N}{w}}
\sum_{s=1}^{r-1}
\binom{w-1}{s-1}\binom{N-w+1}{s}.
\end{equation}
\end{lemma}

\begin{proof}
A binary word of length $N$ with weight $w$ and exactly $s$ runs of ones is described by two independent choices.
First choose the $s$ positive run lengths whose sum is $w$, which gives $\binom{w-1}{s-1}$ possibilities.
Next place the $N-w$ zeros into the $s+1$ zero-gaps around those runs, with the $s-1$ internal gaps required to
receive at least one zero so that adjacent runs remain distinct.
After reserving one zero for each internal gap, the remaining $N-w-(s-1)$ zeros are distributed freely over the
$s+1$ gaps, giving
\[
\binom{N-w-(s-1)+s}{s}
\;=\;
\binom{N-w+1}{s}
\]
choices.
Dividing by the total number $\binom{N}{w}$ of weight-$w$ words yields the stated probability formula,
and summing over $s<r$ gives \eqref{eq:run-lower-tail-exact}.
\end{proof}

\paragraph{Refinements outside the theorem spine.}
Sharper positional refinements for isolated starts, early pairs, multi-episode paths, and paired covers are not used by
the analytic dense+dense theorem below. They are preserved in the scalar dense-inner refinement note together with the
finite scalar diagnostics.

\begin{lemma}[Entropy bound for the run lower tail]
\label{lem:run-lower-tail-entropy}
Fix constants \(0<\omega<1\) and \(0<\rho<\min(\omega,1-\omega)\), and let
\[
w=\lfloor \omega N\rfloor,
\qquad
r=\lfloor \rho N\rfloor.
\]
Assume \(\rho<\omega(1-\omega)\). Then
\begin{equation}
\label{eq:run-lower-tail-entropy}
\PP[R(U)<r]
\;\le\;
N^{O(1)}\,2^{-N\Psi_{\mathrm{run}}(\omega,\rho)},
\end{equation}
where
\begin{equation}
\label{eq:run-rate-function}
\Psi_{\mathrm{run}}(\omega,\rho)
\;:=\;
H_2(\omega)
- \omega H_2\!\Bigl(\frac{\rho}{\omega}\Bigr)
- (1-\omega)H_2\!\Bigl(\frac{\rho}{1-\omega}\Bigr).
\end{equation}
\end{lemma}

\begin{proof}
Insert the exact tail formula \eqref{eq:run-lower-tail-exact}.  For \(s=\lfloor\sigma N\rfloor\), Stirling's formula
bounds each summand by
\[
N^{O(1)}2^{-N\Psi_{\mathrm{run}}(\omega,\sigma)}.
\]
On \(0<\sigma<\omega(1-\omega)\),
\[
\frac{\partial}{\partial\sigma}\Psi_{\mathrm{run}}(\omega,\sigma)
=
\log_2\!\left(\frac{\sigma^2}{(\omega-\sigma)(1-\omega-\sigma)}\right)<0.
\]
Thus \(\Psi_{\mathrm{run}}(\omega,\sigma)\) decreases on the lower-tail range, so the largest summand for
\(s<r=\lfloor\rho N\rfloor\) occurs at the boundary \(\sigma=\rho+O(1/N)\).  Summing at most \(N\) terms gives
\eqref{eq:run-lower-tail-entropy}.  The expanded Stirling derivation is preserved in the scalar dense-inner refinement
note.
\end{proof}

\begin{corollary}[Run lower-tail exponent with \(r\) proportional to \(w\)]
\label{cor:run-lower-tail-alpha}
Fix constants \(0<\omega<1\) and \(0<\alpha<1-\omega\), and let
\[
w=\lfloor \omega N\rfloor,
\qquad
r=\lfloor \alpha w\rfloor.
\]
Then
\begin{equation}
\label{eq:run-lower-tail-alpha}
\PP[R(U)<r]
\;\le\;
N^{O(1)}\,2^{-N\Psi_{\mathrm{run}}^{(w)}(\omega,\alpha)},
\end{equation}
where
\begin{equation}
\label{eq:run-rate-function-alpha}
\Psi_{\mathrm{run}}^{(w)}(\omega,\alpha)
\;:=\;
H_2(\omega)
- \omega H_2(\alpha)
- (1-\omega)H_2\!\Bigl(\frac{\alpha\omega}{1-\omega}\Bigr).
\end{equation}
\end{corollary}

\begin{proof}
Apply Lemma~\ref{lem:run-lower-tail-entropy} with \(\rho=\alpha\omega\).
The condition \(\alpha<1-\omega\) is exactly \(\rho<\omega(1-\omega)\).
\end{proof}

\begin{lemma}[Forced-termination contraction]
\label{lem:dense-forced-termination}
Fix $\delta\in(0,1/2)$ and choose any $\xi>0$. Let
\[
L\ :=\ \left\lceil (2+\xi)\delta N \right\rceil.
\]
Let $U$ be uniform over $\{0,1\}^N$ conditioned on $\wt(U)=w\ge 1$, and let $Y=\Enc_{\mathrm{in}}(U)$ be produced by
\eqref{eq:dense-inner-def}. Then for every integer $r\in\{1,\dots,w\}$,
\begin{equation}
\label{eq:dense-forced-termination}
\PP[\wt(Y)\le \delta N \mid \wt(U)=w]
\ \le\
\PP[R(U)<r]\ +\ (L2^{-m})^{r}\ +\ r\cdot \PP[\mathrm{Binomial}(L,\tfrac12)\le \delta N].
\end{equation}
Moreover,
\begin{equation}
\label{eq:runs-adjacent-bound}
\PP[R(U)<w]\ \le\ \frac{w(w-1)}{N-1}.
\end{equation}
In particular, for $w\ll \sqrt{N}$ (e.g.\ $w=O(\log N)$), taking $r=w$ yields a geometric-in-$w$ term
$(L2^{-m})^w$ up to a negligible $O(w^2/N)$ correction.
\end{lemma}

\begin{proof}
\emph{Step 1 (runs initiate ON-episodes).}
When the encoder is OFF ($s_t=0$), we have $y_t=u_t$.
Therefore, at the first position of any run of ones in $u$, the encoder outputs $1$, and the next state becomes nonzero,
so the encoder enters ON immediately after the run begins.
Thus each run of ones initiates an ON-episode.

\emph{Step 2 (small output weight forces fast termination or a binomial deviation).}
Condition on the event $\{R(U)\ge r\}$ and select any $r$ distinct runs.
Consider the corresponding $r$ ON-episodes (they occur on disjoint time intervals since the runs are disjoint).

If any selected ON-episode has length at least $L$, then on the sub-interval where the encoder stays ON,
Lemma~\ref{lem:dense-on-time-domination} gives fair-coin domination.
Since $L=(2+\xi)\delta N>2\delta N$, the probability that this single episode contributes at most $\delta N$ ones is
upper bounded by $\PP[\mathrm{Binomial}(L,1/2)\le \delta N]$ (monotonicity).
By a union bound over the $r$ selected episodes, the probability that some selected episode is long and still contributes
$\le \delta N$ ones is at most
\[
r\cdot \PP[\mathrm{Binomial}(L,1/2)\le \delta N].
\]

Otherwise, all $r$ selected ON-episodes have length $<L$.
An ON-episode ends only via an ON$\to$OFF transition, which by Lemma~\ref{lem:dense-off-requires-mzeros} requires a
length-$m$ all-zero window in the output while ON.
For a given ON-episode to terminate within $L$ steps, there are at most $L$ candidate start indices for such a window;
for each fixed candidate, Lemma~\ref{lem:dense-on-time-domination} gives probability at most \(2^{-m}\) for an
all-zero window.
Thus, by a union bound, a fixed ON-episode terminates within $L$ steps with probability at most $L2^{-m}$.
Since tap vectors are independent over time, the termination events for the $r$ disjoint ON-episodes are independent
conditional on the input pattern. Hence the probability that all $r$ selected ON-episodes terminate within $L$ is
at most $(L2^{-m})^r$.

Finally, union bound over the complementary event $\{R(U)<r\}$ yields \eqref{eq:dense-forced-termination}.

\emph{Step 3 (adjacent-ones bound).}
We have $R(U)<w$ iff $U$ contains an adjacent pair of ones.
Let
\[
A:=\sum_{i=1}^{N-1}\mathbf{1}\{U_i=U_{i+1}=1\}.
\]
Then $\{R(U)<w\}\subseteq\{A\ge 1\}$, so $\PP[R(U)<w]\le \EE[A]$.
For uniform weight-$w$ vectors,
\[
\PP[U_i=U_{i+1}=1]=\frac{\binom{N-2}{w-2}}{\binom{N}{w}}=\frac{w(w-1)}{N(N-1)}.
\]
Therefore $\EE[A]=(N-1)\cdot \frac{w(w-1)}{N(N-1)}=\frac{w(w-1)}{N}$, which implies \eqref{eq:runs-adjacent-bound}
(up to the minor $N$ vs.\ $N-1$ normalization).
\end{proof}

\paragraph{Linear OFF-budget refinement.}
A sharper linear-weight route can try to bound the total post-start OFF time directly by separating long zero gaps from
short-gap witness packings. That proof target is also preserved in
the scalar dense-inner refinement note; the main theorem below only uses the run-tail and forced-termination envelopes.

% ============================================================
\subsection{Inner interface theorem (what the outer+interleaver proof consumes)}
\label{subsec:inner-dense-interface}

We package the previous lemmas into an explicit interface for Section~\ref{sec:framework}.

\begin{theorem}[Dense inner interface for the framework]
\label{thm:inner-dense-interface}
Let $\Ctwo$ be the dense time-varying recursive inner encoder \eqref{eq:dense-inner-def} with memory $m$.
Fix $\delta\in(0,1/2)$.
For each $w\ge 1$, let $U$ be uniform over $\{0,1\}^N$ conditioned on $\wt(U)=w$ and let $Y=\Enc_{\mathrm{in}}(U)$.

Then:
\begin{enumerate}
\item (\textbf{General tail, all $w$}) For every $\varepsilon\in(0,1)$,
\[
p_w(\delta)\ :=\ \PP[\wt(Y)\le \delta N\mid \wt(U)=w]
\ \le\
(1-\varepsilon)^w\ +\ N2^{-m}\ +\ q_{\delta,\varepsilon}(N),
\]
with $q_{\delta,\varepsilon}(N)$ as in Lemma~\ref{lem:dense-three-term-tail}.
\item (\textbf{Forced-termination contraction, tunable in $r$}) For every $\xi>0$ and
$L=\lceil(2+\xi)\delta N\rceil$, and for every $r\in\{1,\dots,w\}$,
\[
p_w(\delta)
\ \le\
\PP[R(U)<r]\ +\ (L2^{-m})^{r}\ +\ r\cdot \PP[\mathrm{Binomial}(L,\tfrac12)\le \delta N].
\]
\end{enumerate}
\end{theorem}

\begin{corollary}[Optimization-ready dense-inner envelopes]
\label{cor:dense-inner-piecewise}
Fix $\delta\in(0,1/2)$.
Choose any $\varepsilon\in(0,1)$ and any $\xi>0$, and define
\[
L:=\lceil(2+\xi)\delta N\rceil,
\]
\[
q_{\delta,\varepsilon}(N)
:=
\PP\!\left[\ \Binomial(\lfloor(1-\varepsilon)N\rfloor,\tfrac12)\le \delta N\ \right],
\]
\[
b_{\delta,\xi}(N)
:=
\PP\!\left[\ \Binomial(L,\tfrac12)\le \delta N\ \right].
\]
For $1\le r\le w$, define the exact run lower-tail
\[
\mathrm{RunTail}_N(w,r)
\;:=\;
\frac{1}{\binom{N}{w}}
\sum_{s=1}^{r-1}
\binom{w-1}{s-1}\binom{N-w+1}{s}.
\]
Then for every $w\ge 1$,
\begin{equation}
\label{eq:dense-inner-global-envelope}
p_w(\delta)
\;\le\;
(1-\varepsilon)^w + N2^{-m} + q_{\delta,\varepsilon}(N),
\end{equation}
and for every $r\in\{1,\dots,w\}$,
\begin{equation}
\label{eq:dense-inner-run-envelope}
p_w(\delta)
\;\le\;
\mathrm{RunTail}_N(w,r) + (L2^{-m})^r + r\,b_{\delta,\xi}(N).
\end{equation}
In particular,
\begin{equation}
\label{eq:dense-inner-run-opt-envelope}
p_w(\delta)
\;\le\;
B_w^{\mathrm{run}}(\delta;\xi,m)
\;:=\;
\min_{1\le r\le w}
\left(
\mathrm{RunTail}_N(w,r) + (L2^{-m})^r + r\,b_{\delta,\xi}(N)
\right),
\end{equation}
and also
\begin{equation}
\label{eq:dense-inner-smallw-envelope}
p_w(\delta)
\;\le\;
\frac{w(w-1)}{N-1}
\;+\;
\frac{\binom{L-w+1}{w}}{\binom{N-w+1}{w}}
\;+\;
w2^{-(m-1)}
\;+\;
b_{\delta,\xi}(N).
\end{equation}
Consequently,
\begin{equation}
\label{eq:dense-inner-best-envelope}
p_w(\delta)
\;\le\;
\min\!\left\{
B_w^{\mathrm{run}}(\delta;\xi,m),
(1-\varepsilon)^w + N2^{-m} + q_{\delta,\varepsilon}(N)
\right\}.
\end{equation}
\end{corollary}

\begin{proof}
The global bound \eqref{eq:dense-inner-global-envelope} is exactly the first part of
Theorem~\ref{thm:inner-dense-interface}.
The bound \eqref{eq:dense-inner-run-envelope} is the second part of
Theorem~\ref{thm:inner-dense-interface}, together with the exact run-tail identity
\eqref{eq:run-lower-tail-exact}.
The optimized form \eqref{eq:dense-inner-run-opt-envelope} is immediate by minimizing over $r$.
For the specialization \eqref{eq:dense-inner-smallw-envelope}, split first according to whether \(R(U)=w\). The
non-isolated part is bounded by \(\PP[R(U)<w]\le w(w-1)/(N-1)\) from \eqref{eq:runs-adjacent-bound}. On the isolated
slice, write the one positions as \(X_1<\cdots<X_w\) and shift to \(Y_i=X_i-i+1\). Then \(Y_1<\cdots<Y_w\) is a
uniform \(w\)-subset of \([N-w+1]\), so the probability that the first input one lies in the final \(L\) coordinates is
\[
\frac{\binom{L-w+1}{w}}{\binom{N-w+1}{w}}.
\]
On the complementary event, the first ON episode starts by time \(N-L\). If it survives for \(L\) output times, the
low-weight event is bounded by \(b_{\delta,\xi}(N)\) via Lemma~\ref{lem:fixedtap-surviving-episode-tail}. If it
terminates within \(L\) steps, Lemma~\ref{lem:fixedtap-one-episode-termination} bounds the probability by
\(B_{T_1}(U;L)2^{-(m-1)}\le w2^{-(m-1)}\). Summing these cases gives
\eqref{eq:dense-inner-smallw-envelope}.
The best-envelope form \eqref{eq:dense-inner-best-envelope} is then immediate.
\end{proof}

\paragraph{Why these are the right envelopes to optimize.}
The bound \eqref{eq:dense-inner-global-envelope} is designed for the linear-weight regime:
its first term decays exponentially in $w$, its second term exposes the ON$\to$OFF memory penalty $N2^{-m}$,
and its third term is a pure large-deviation term in $N$.
The run-aware bound \eqref{eq:dense-inner-run-opt-envelope} is the legacy fully-random-tap tool for finite-length
comparison:
it keeps the exact combinatorics of the interleaved run count,
it turns the OFF mechanism into the genuinely multiplicative term $(L2^{-m})^r$,
and it can be optimized over $r$ separately for each weight $w$.
The simpler bound \eqref{eq:dense-inner-smallw-envelope} is a convenient specialization for
the fixed-tap regime where random interleaving typically creates nearly \(w\) isolated starts and genuine termination
candidates must be aligned with later input-one positions.

\begin{theorem}[Dense-inner linear-weight exponent]
\label{prop:dense-inner-linear-exponent}
Fix constants $\delta\in(0,1/2)$, $\xi>0$, and $\theta\in(0,1)$.
Let
\[
L:=\min\{N,\lceil(2+\xi)\delta N\rceil\},
\qquad
m=\gamma\log_2 N
\]
with $\gamma>1$, and define
\[
\lambda_{\delta,\xi}
:=
\min\{(2+\xi)\delta,1\},
\]
and
\begin{equation}
\label{eq:dense-inner-binomial-exponent}
E_{\mathrm{bin}}(\delta,\xi)
\;:=\;
\lambda_{\delta,\xi}
\left(1-H_2\!\Bigl(\frac{\delta}{\lambda_{\delta,\xi}}\Bigr)\right).
\end{equation}
For a fixed weight fraction $\omega\in(0,1)$ satisfying $\theta<1-\omega$, let
\[
w=\lfloor \omega N\rfloor,
\qquad
r=\lfloor \theta w\rfloor.
\]
Then the dense inner satisfies
\begin{equation}
\label{eq:dense-inner-linear-exponent}
p_w(\delta)
\;\le\;
N^{O(1)}\,
2^{-N E_{\mathrm{in}}(\omega;\theta,\delta,\xi)},
\end{equation}
where
\begin{equation}
\label{eq:dense-inner-linear-rate-function}
E_{\mathrm{in}}(\omega;\theta,\delta,\xi)
\;:=\;
\min\!\left\{
\Psi_{\mathrm{run}}^{(w)}(\omega,\theta),\
E_{\mathrm{bin}}(\delta,\xi)
\right\},
\end{equation}
and $\Psi_{\mathrm{run}}^{(w)}(\omega,\theta)$ is defined in
\eqref{eq:run-rate-function-alpha}.
\end{theorem}

\begin{proof}
Apply \eqref{eq:dense-inner-run-envelope} with the choice $r=\lfloor \theta w\rfloor$.
The run-deficit term is controlled by Corollary~\ref{cor:run-lower-tail-alpha}:
\[
\mathrm{RunTail}_N(w,r)
\;\le\;
N^{O(1)}\,2^{-N\Psi_{\mathrm{run}}^{(w)}(\omega,\theta)}.
\]

Next,
\[
L2^{-m}
\;=\;
\Theta(N)\cdot N^{-\gamma}
\;=\;
N^{1-\gamma+o(1)},
\]
so since $r=\Theta(N)$ and $\gamma>1$,
\[
(L2^{-m})^r
\;=\;
2^{-\Omega(N\log N)}.
\]

For the binomial term, note that
\[
\frac{\delta N}{L}
\;=\;
\frac{\delta}{\lambda_{\delta,\xi}}+O\!\left(\frac{1}{N}\right).
\]
Since \(\delta/\lambda_{\delta,\xi}<1/2\), a standard Stirling/Chernoff bound gives
\[
\PP[\Binomial(L,\tfrac12)\le \delta N]
\;\le\;
N^{O(1)}\,2^{-L\left(1-H_2(\delta/\lambda_{\delta,\xi})\right)}
\;\le\;
N^{O(1)}\,2^{-N E_{\mathrm{bin}}(\delta,\xi)}.
\]
Multiplying by $r=O(N)$ only changes the polynomial prefactor.

Combining the three terms in \eqref{eq:dense-inner-run-envelope} yields
\eqref{eq:dense-inner-linear-exponent}.
\end{proof}




\begin{corollary}[Geometric dense-inner envelope above a logarithmic cutoff]
\label{cor:dense-inner-geometric-cutoff}
Consider the dense time-varying recursive inner encoder with memory
\[
m \;:=\; \gamma \log_2 n
\qquad\text{for a constant }\gamma>1.
\]
Fix $\delta\in(0,1/2)$ and choose any fixed $\varepsilon\in(0,1)$.
Let $p_w(\delta)$ denote the slice-to-tail probability from \eqref{eq:dense-inner-global-envelope}.
Assume the dense-inner tail bound (proved earlier in this section) gives
\begin{equation}
\label{eq:dense-three-term-recall}
p_w(\delta)
\;\le\;
(1-\varepsilon)^w \;+\; n\,2^{-m} \;+\; q_{\delta,\varepsilon}(n),
\end{equation}
where $q_{\delta,\varepsilon}(n)\le 2^{-c n}$ for some $c=c(\delta,\varepsilon)>0$.

Then there exist constants $\rho\in(0,1)$, $a\ge 0$, and $\kappa>0$
(depending only on $\delta,\varepsilon,\gamma$) such that
for every fixed cutoff $w_0(n)\ge \kappa\log_2 n$,
for all sufficiently large $n$, and for all $w\ge w_0(n)$,
\[
p_w(\delta)\;\le\; n^{a}\,\rho^{\,w}.
\]
\end{corollary}

\begin{proof}
Let $\rho:=1-\varepsilon\in(0,1)$.
We bound each term in \eqref{eq:dense-three-term-recall} by a common expression $n^a\rho^w$ for all
$w\ge w_0(n)\ge \kappa\log_2 n$.

\paragraph{Term 1.}
$(1-\varepsilon)^w=\rho^w\le n^a\rho^w$ for any $a\ge 0$.

\paragraph{Term 2.}
Since $m=\gamma\log_2 n$, we have $n\,2^{-m}=n\cdot n^{-\gamma}=n^{1-\gamma}$.
Then for all $w\ge w_0(n)\ge \kappa\log_2 n$,
\[
\rho^{\,w}\ \le\ \rho^{\,\kappa\log_2 n}
\ =\ n^{-\kappa\log_2(1/\rho)}.
\]
Choose any constant $a$ such that
\[
a - \kappa\log_2(1/\rho)\ \ge\ 1-\gamma.
\]
Then $n^{1-\gamma}\le n^{a}\rho^{\,w}$ for all $w\ge w_0(n)$.

\paragraph{Term 3.}
Since $q_{\delta,\varepsilon}(n)\le 2^{-cn}$ and $\rho^w\ge \rho^n$ is exponentially small in $n$ as well,
we have $q_{\delta,\varepsilon}(n)\le n^{a}\rho^{\,w}$ for all sufficiently large $n$ after increasing $a$ by an absolute constant if needed.

Combining the three bounds yields $p_w(\delta)\le n^a\rho^w$ for all $w\ge w_0(n)$.
\end{proof}

\begin{corollary}[Constant-factor geometric envelope]
\label{cor:dense-inner-constant-geometric}
Fix \(\delta\in(0,1/2)\) and choose \(\rho\in(0,1)\) with \(2\delta<\rho\).
Let \(m=\gamma\log_2 n\), and fix any logarithmic cutoff \(w_0(n)\ge \kappa\log_2 n\).
If
\[
\gamma \;>\; 1+\kappa\log_2\!\Bigl(\frac{1}{\rho}\Bigr),
\]
then for all sufficiently large \(n\) and all \(w\ge w_0(n)\),
\[
p_w(\delta)\le 3\rho^w.
\]
\end{corollary}

\begin{proof}
Apply Lemma~\ref{lem:dense-three-term-tail} with \(\varepsilon=1-\rho\).
The first term is exactly \(\rho^w\).
Since \(2\delta<\rho\), the binomial tail
\[
q_{\delta,1-\rho}(n)
=
\PP\!\left[\Binomial(\lfloor \rho n\rfloor,\tfrac12)\le \delta n\right]
\]
is exponentially small in \(n\), hence \(q_{\delta,1-\rho}(n)\le \rho^w\) for all sufficiently large \(n\) and all
\(w\ge \kappa\log_2 n\).
Also,
\[
n2^{-m}=n^{1-\gamma}
\le
n^{-\kappa\log_2(1/\rho)}
\le
\rho^w
\]
for all \(w\ge w_0(n)\), where the first inequality uses
\(\gamma>1+\kappa\log_2(1/\rho)\).
Summing the three terms proves the claim.
\end{proof}

\begin{remark}[The ON$\to$OFF term and why it is harmless in the framework]
\label{rem:dense-off-harmless}
The term $n2^{-m}$ in \eqref{eq:dense-three-term-recall} corresponds to the event that, while ON, the output stream
contains a length-$m$ all-zero window, causing an ON$\to$OFF transition.
For dense taps the expected ON lifetime under zero input is $\Theta(2^m)$, hence this term is unavoidable in a pointwise bound.
However, the serial concatenation framework never uses $p_w(\delta)$ for $w<d_{\min}(\Cone)$,
and for $w\ge d_{\min}(\Cone)=\Omega(\log n)$ the geometric factor $\rho^w$ yields an additional $n^{-\Omega(1)}$
that easily absorbs $n2^{-m}=n^{1-\gamma}$ as soon as $\gamma>1$.
Thus the dense inner conforms to the single inner interface without weakening the end-to-end guarantee.
\end{remark}




\paragraph{Work complexity (why this is the ``dense'' baseline).}
Each step computes $\langle k_t,s_t\rangle$ with a \emph{dense} $m$-bit mask, i.e.\ $\Theta(m)$ taps per symbol.
Thus total inner work is $\Theta(Nm)=\Theta(N\log N)$ for $m=\gamma\log N$.

\paragraph{Parameter-use summary.}
\begin{itemize}
\item The term $(1-\varepsilon)^w$ is the clean exponential-in-$w$ contraction for intermediate weights.
\item The term $N2^{-m}$ is the rare ON$\to$OFF event. In tight finite-$N$ bounds, it is better to \emph{budget}
its cumulative contribution in the weight sum (rather than carry it pointwise for every $w$).
\item The binomial tail $q_{\delta,\varepsilon}(N)$ provides the exponentially small-in-$N$ large-deviation behavior
needed for the linear-weight regime when union bounding over exponentially many outer codewords.
\item Lemma~\ref{lem:dense-forced-termination} replaces the pointwise $N2^{-m}$ by $(L2^{-m})^r$ for small $w$,
using the fact that a random interleaver typically spreads $w$ ones into $\approx w$ runs when $w\ll \sqrt N$.
\end{itemize}

% ============================================================
